\chapter{Partial Fractions}

Decompose each rational expression. Problems 2 and 3 are \emph{improper}
(numerator degree $\ge$ denominator degree): perform polynomial division
first.

\section{Core Rules}

\begin{rulebox}{What the method does, and why}
A partial fraction decomposition rewrites one complicated fraction as a sum
of simple ones:
\[
  \frac{7x+1}{(x+2)(x-3)}=\frac{A}{x+2}+\frac{B}{x-3} .
\]
It is the reverse of putting fractions over a common denominator. The point
is that each piece on the right is something you can integrate, invert, or
sum, while the left side is not.
\end{rulebox}

\begin{rulebox}{Step 0, which must not be skipped: make it proper}
The method applies only when $\deg(\text{numerator})<\deg(\text{denominator})$.
If not, \textbf{divide first}:
\[
  \frac{x^{2}}{(x-1)(x-2)}
  =\frac{x^{2}}{x^{2}-3x+2}
  =1+\frac{3x-2}{x^{2}-3x+2} ,
\]
and decompose only the remainder. An improper fraction fed straight into the
ansatz produces an unsolvable or inconsistent system.
\end{rulebox}

\begin{rulebox}{Choosing the ansatz}
Factor the denominator completely, then contribute one term per factor:
\begin{itemize}[leftmargin=1.6em,itemsep=0.3ex]
  \item \textbf{Distinct linear factor} $(x-a)$: one term $\dfrac{A}{x-a}$.
  \item \textbf{Repeated linear factor} $(x-a)^{k}$: \textbf{all} $k$ terms,
        \[
          \frac{A_{1}}{x-a}+\frac{A_{2}}{(x-a)^{2}}+\cdots+\frac{A_{k}}{(x-a)^{k}} .
        \]
  \item \textbf{Irreducible quadratic} $(x^{2}+px+q)$ with $p^{2}<4q$:
        a \emph{linear} numerator, $\dfrac{Bx+C}{x^{2}+px+q}$.
\end{itemize}
The numerator always has degree one less than its denominator factor.
\end{rulebox}

\begin{rulebox}{Finding the constants: two techniques}
Multiply through by the full denominator to get a polynomial identity, then
either
\begin{itemize}[leftmargin=1.6em,itemsep=0.3ex]
  \item \textbf{substitute the roots} (the ``cover-up'' method) --- each root
        of a linear factor eliminates every term but one, giving a constant
        immediately; or
  \item \textbf{compare coefficients} of like powers, giving a small linear
        system.
\end{itemize}
Roots are faster and handle the simple constants; coefficient comparison
supplies the ones the roots cannot reach, which is exactly the situation with
a repeated factor.
\end{rulebox}

\section{Worked Examples}

\begin{worked}{(a) Decompose $\dfrac{x^{2}-5}{x(x+1)^{2}}$}
\wstep{Step 1. Check properness, then write the ansatz.}
The numerator has degree $2$, the denominator degree $3$: proper, so no
division is needed. The factor $(x+1)$ is repeated, so it contributes
\emph{two} terms:
\[
  \frac{x^{2}-5}{x(x+1)^{2}}
  =\frac{A}{x}+\frac{B}{x+1}+\frac{C}{(x+1)^{2}} .
\]
\wtrap{Writing only $\dfrac{A}{x}+\dfrac{C}{(x+1)^{2}}$ is the classic error.
Two unknowns cannot match three coefficients, and the resulting system will
be inconsistent --- which at least tells you something is wrong. A repeated
factor of multiplicity $k$ needs all $k$ powers.}

\wstep{Step 2. Clear the denominators.}
Multiplying both sides by $x(x+1)^{2}$:
\[
  x^{2}-5=A(x+1)^{2}+Bx(x+1)+Cx .
\]
This is an identity: it holds for \emph{every} $x$, which is what licenses
the substitutions that follow.

\wstep{Step 3. Substitute the roots to find the easy constants.}
\[
  x=0:\quad -5=A(1)+0+0 \;\Longrightarrow\; A=-5 .
\]
\[
  x=-1:\quad 1-5=0+0+C(-1) \;\Longrightarrow\; -4=-C \;\Longrightarrow\; C=4 .
\]

\wstep{Step 4. Get $B$ by comparing coefficients.}
No substitution isolates $B$ --- that is characteristic of a repeated
factor. Compare the coefficient of $x^{2}$ on both sides:
\[
  1=A+B \;\Longrightarrow\; B=1-(-5)=6 .
\]
\wtrap{Only $A(x+1)^{2}$ and $Bx(x+1)$ produce an $x^{2}$; $Cx$ does not.
Expand mentally and take just the one coefficient you need rather than
multiplying everything out.}

\wstep{Step 5. Verify at a value not yet used.}
\[
  \frac{x^{2}-5}{x(x+1)^{2}}=\boxed{-\frac{5}{x}+\frac{6}{x+1}+\frac{4}{(x+1)^{2}}}
\]
Test $x=1$: the left side is $(1-5)/(1\cdot4)=-1$; the right side is
$-5+3+1=-1$. \checkmark\ One substitution catches nearly every arithmetic
slip, and it costs ten seconds.
\end{worked}

\begin{worked}{(b) The improper case: $\dfrac{x^{2}}{(x-1)(x-2)}$}
\wstep{Step 1. Compare degrees --- and stop.}
Numerator degree $2$, denominator degree $2$. The fraction is improper, so
the ansatz $\frac{A}{x-1}+\frac{B}{x-2}$ cannot possibly work: that sum
tends to $0$ as $x\to\infty$, while the left side tends to $1$.

\wstep{Step 2. Divide, by adding and subtracting rather than by long division.}
The denominator expands to $x^{2}-3x+2$, so
\[
  x^{2}=(x^{2}-3x+2)+(3x-2)
  \quad\Longrightarrow\quad
  \frac{x^{2}}{(x-1)(x-2)}=1+\frac{3x-2}{(x-1)(x-2)} .
\]

\wstep{Step 3. Decompose the proper remainder.}
\[
  \frac{3x-2}{(x-1)(x-2)}=\frac{A}{x-1}+\frac{B}{x-2}
  \quad\Longrightarrow\quad
  3x-2=A(x-2)+B(x-1).
\]
\[
  x=1:\ 1=A(-1)\Rightarrow A=-1 ;\qquad
  x=2:\ 4=B(1)\Rightarrow B=4 .
\]

\wstep{Step 4. Reassemble --- and do not lose the $1$.}
\[
  \frac{x^{2}}{(x-1)(x-2)}=\boxed{1-\frac{1}{x-1}+\frac{4}{x-2}}
\]
Check at $x=0$: left side $0$; right side $1-\frac{1}{-1}+\frac{4}{-2}
=1+1-2=0$. \checkmark
\wtrap{The polynomial part is part of the answer. Decomposing the remainder
and then reporting only the fractions loses a whole term --- and, in the
integration volume, loses the $x$ in the antiderivative.}
\end{worked}

\section{Practice Problems}

\begin{problems}
  \item $\displaystyle\frac{7x+1}{(x+2)(x-3)}$
  \item $\displaystyle\frac{x^{2}}{(x-1)(x-2)}$ (improper)
  \item $\displaystyle\frac{2x^{2}+5x-3}{(x+1)(x-2)}$ (improper)
  \item $\displaystyle\frac{4x-5}{x^{2}-3x+2}$
  \item $\displaystyle\frac{6x+1}{(x-1)(x+3)}$
  \item $\displaystyle\frac{x-5}{x(x-4)}$
  \item $\displaystyle\frac{1}{(x+4)(x+5)}$
  \item $\displaystyle\frac{x-4}{x^{2}+2x-8}$
  \item $\displaystyle\frac{1}{x^{2}-7x+12}$
  \item $\displaystyle\frac{2x^{2}+3}{x(x-1)(x+2)}$
  \item $\displaystyle\frac{x+2}{(x-1)^{2}}$
  \item $\displaystyle\frac{5x}{(x+2)^{2}}$
  \item $\displaystyle\frac{1}{x^{2}(x+1)}$
  \item $\displaystyle\frac{x^{2}-5}{x(x+1)^{2}}$
  \item $\displaystyle\frac{2x+3}{x^{2}+4x}$
  \item $\displaystyle\frac{1}{(x-1)(x^{2}-4)}$
  \item $\displaystyle\frac{5-x}{(x+3)(x-1)}$
  \item $\displaystyle\frac{4x-7}{(x-2)(x+1)(x+3)}$
  \item $\displaystyle\frac{1}{(s+a)(s+b)}$ ($a\neq b$, symbolic)
  \item $\displaystyle\frac{2K-1}{K(1-K)}$ (logistic-style)
  \item $\displaystyle\frac{3x^{2}+4x+1}{x^{3}}$
  \item $\displaystyle\frac{1}{(2s+1)(s-2)}$ (Laplace style)
  \item $\displaystyle\frac{6}{x^{2}-25}$
  \item $\displaystyle\frac{\mathrm{e}^{x}}{\mathrm{e}^{2x}-1}$
        (substitute $u=\mathrm{e}^{x}$ first)
  \item $\displaystyle\frac{3t+1}{t(t-1)(t-2)}$
  \item $\displaystyle\frac{x+4}{x^{3}+x^{2}}$
  \item $\displaystyle\frac{2y-3}{(y+1)^{2}}$
  \item $\displaystyle\frac{4x-5}{(x-3)^{2}}$
  \item $\displaystyle\frac{1}{(x+a)(x+b)}$ ($a\neq b$, symbolic)
  \item Telescoping: write $\displaystyle\frac{1}{(k-1)k}$ ($k\ge 2$) as a
        difference of two terms.
\end{problems}

\section{Solutions}

\begin{solutions}
  \item $\displaystyle\frac{13/5}{x+2}+\frac{22/5}{x-3}$
  \item $1-\displaystyle\frac{1}{x-1}+\frac{4}{x-2}$
  \item $2+\displaystyle\frac{2}{x+1}+\frac{5}{x-2}$
  \item $\displaystyle\frac{1}{x-1}+\frac{3}{x-2}$
  \item $\displaystyle\frac{7/4}{x-1}+\frac{17/4}{x+3}$
  \item $\displaystyle\frac{5/4}{x}-\frac{1/4}{x-4}$
  \item $\displaystyle\frac{1}{x+4}-\frac{1}{x+5}$
  \item $\displaystyle-\frac{1/3}{x-2}+\frac{4/3}{x+4}$
  \item $\displaystyle-\frac{1}{x-3}+\frac{1}{x-4}$
  \item $\displaystyle-\frac{3/2}{x}+\frac{5/3}{x-1}+\frac{11/6}{x+2}$
  \item $\displaystyle\frac{1}{x-1}+\frac{3}{(x-1)^{2}}$
  \item $\displaystyle\frac{5}{x+2}-\frac{10}{(x+2)^{2}}$
  \item $\displaystyle-\frac{1}{x}+\frac{1}{x^{2}}+\frac{1}{x+1}$
  \item $\displaystyle-\frac{5}{x}+\frac{6}{x+1}+\frac{4}{(x+1)^{2}}$
  \item $\displaystyle\frac{3/4}{x}+\frac{5/4}{x+4}$
  \item $\displaystyle-\frac{1/3}{x-1}+\frac{1/4}{x-2}+\frac{1/12}{x+2}$
  \item $\displaystyle-\frac{2}{x+3}+\frac{1}{x-1}$
  \item $\displaystyle\frac{1/15}{x-2}+\frac{11/6}{x+1}-\frac{19/10}{x+3}$
  \item $\displaystyle\frac{1}{b-a}\left(\frac{1}{s+a}-\frac{1}{s+b}\right)$
  \item $\displaystyle-\frac{1}{K}+\frac{1}{1-K}$
  \item $\displaystyle\frac{3}{x}+\frac{4}{x^{2}}+\frac{1}{x^{3}}$
  \item $\displaystyle-\frac{2/5}{2s+1}+\frac{1/5}{s-2}$
  \item $\displaystyle\frac{3/5}{x-5}-\frac{3/5}{x+5}$
  \item $\displaystyle\frac{1/2}{\mathrm{e}^{x}-1}+\frac{1/2}{\mathrm{e}^{x}+1}$
  \item $\displaystyle\frac{1/2}{t}-\frac{4}{t-1}+\frac{7/2}{t-2}$
  \item $\displaystyle-\frac{3}{x}+\frac{4}{x^{2}}+\frac{3}{x+1}$
  \item $\displaystyle\frac{2}{y+1}-\frac{5}{(y+1)^{2}}$
  \item $\displaystyle\frac{4}{x-3}+\frac{7}{(x-3)^{2}}$
  \item $\displaystyle\frac{1}{b-a}\left(\frac{1}{x+a}-\frac{1}{x+b}\right)$
  \item $\displaystyle\frac{1}{k-1}-\frac{1}{k}$
\end{solutions}
